\documentclass{article}

% if you need to pass options to natbib, use, e.g.:
%     \PassOptionsToPackage{numbers, compress}{natbib}
% before loading neurips_2024

% use this option for ANONYMOUS submission
 \usepackage{neurips_2024}

% use this option for NON-ANONYMOUS compiling. Not allowed for submission.
% \usepackage[final]{neurips_2024}

\usepackage[utf8]{inputenc}
\usepackage{natbib}
\usepackage[T1]{fontenc}
\usepackage{hyperref}
\usepackage{url}
\usepackage{booktabs}
\usepackage{amsfonts}
\usepackage{nicefrac}
\usepackage{microtype}
\usepackage{xcolor}
\usepackage{graphicx}
\usepackage{listings}
\usepackage{float}
\usepackage{amsmath}

% Code listing style
\lstset{
    language=Python,
    basicstyle=\ttfamily\small,
    keywordstyle=\color{blue},
    stringstyle=\color{red},
    commentstyle=\color{green!60!black},
    numbers=none,
    frame=single,
    breaklines=true,
    showstringspaces=false
}

\title{Vision Transformer: An Image is Worth 16x16 Words}

\author{%
  Rachana Vijay \\
  Illinois Institute of Technology\\
}

\begin{document}

\maketitle

\begin{abstract}
This project presents a CPU-friendly implementation and analysis of the Vision Transformer (ViT), a transformer-based architecture for image classification. We adapt the transformer architecture, originally designed for natural language processing, to computer vision tasks by treating image patches as sequences. Our miniaturized ViT model (36,746 parameters) is trained on a subset of FashionMNIST (1,500 training, 500 test samples), achieving 61.02\% validation accuracy. The implementation demonstrates core ViT principles including patch embedding, multi-head self-attention, and transformer encoder layers. Key insights include: (1) grayscale FashionMNIST simplifies attention pattern interpretation, (2) structured clothing shapes are well-suited for global attention mechanisms, and (3) miniaturized ViTs can train efficiently on standard laptops. We provide comprehensive analysis of the architecture, training process, and discuss strengths and limitations of this approach for minimal CPU-ready working examples.
\end{abstract}

\section{Background}

\subsection{Motivation}

Traditional convolutional neural networks (CNNs) have dominated computer vision tasks for decades, relying on inductive biases such as translation equivariance and locality. Although effective, CNNs require extensive architectural engineering and may not fully exploit global relationships in images. The Vision Transformer paper \citep{dosovitskiy2021vit} addresses this limitation by proposing a pure transformer architecture for image classification, eliminating the need for convolution operations entirely.

The key problem addressed is whether the self-attention mechanism, which has revolutionized natural language processing, can be effectively applied to image classification without relying on CNN-specific inductive biases. This question is fundamental as it challenges long-standing assumptions about how visual information should be processed.

\subsection{Key Breakthrough}

The Vision Transformer achieves its breakthrough through several key innovations:

\textbf{1. Patch-based Image Representation:} Instead of processing images pixel-by-pixel, ViT divides images into fixed-size patches (e.g., $16 \times 16$ pixels) and treats each patch as a token, similar to words in NLP. This transformation allows the direct application of transformer architectures to images.

\textbf{2. Learnable Positional Embeddings:} Since transformers are permutation-invariant, ViT adds learnable positional embeddings to patch embeddings to preserve spatial information. This is crucial for understanding image structure.

\textbf{3. Class Token:} Following BERT \citep{devlin2018bert}, ViT introduces a learnable classification token (CLS token) that aggregates information from all patches through self-attention, enabling global image understanding.

\textbf{4. Pre-training on Large Datasets:} The paper demonstrates that when pre-trained on large datasets (e.g., ImageNet-21k with 14M images), ViT can match or exceed CNN performance while requiring fewer computational resources for fine-tuning.

These innovations enable transformers to process images as sequences, leveraging the same powerful self-attention mechanism that transformed NLP \citep{vaswani2017attention}, without requiring convolution-specific architectural constraints.

\section{Model Architecture}

\subsection{Analysis: Part-by-Part}

The Vision Transformer consists of several key components, each serving a specific purpose in the image classification pipeline.

\subsubsection{Patch Embedding}

The first component converts images into a sequence of patch embeddings. For an image $\mathbf{x} \in \mathbb{R}^{H \times W \times C}$ divided into patches of size $P \times P$, we obtain $N = HW/P^2$ patches. Each patch $\mathbf{x}_p^i \in \mathbb{R}^{P^2 \cdot C}$ is flattened and linearly projected into a $D$-dimensional embedding space using a learnable projection matrix $\mathbf{E} \in \mathbb{R}^{(P^2 \cdot C) \times D}$:

\begin{equation}
\mathbf{z}_0 = [\mathbf{x}_{\text{class}}; \mathbf{x}_p^1\mathbf{E}; \mathbf{x}_p^2\mathbf{E}; \cdots; \mathbf{x}_p^N\mathbf{E}] + \mathbf{E}_{\text{pos}}
\label{eq:patch_embedding}
\end{equation}

where $\mathbf{x}_{\text{class}} \in \mathbb{R}^D$ is the learnable class token, and $\mathbf{E}_{\text{pos}} \in \mathbb{R}^{(N+1) \times D}$ represents learnable positional embeddings that encode spatial information.

\begin{lstlisting}[caption={Patch Embedding Implementation}]
class PatchEmbedding(nn.Module):
    def __init__(self, in_channels=3, patch_size=4, 
                 emb_size=64, img_size=32):
        super().__init__()
        self.patch_size = patch_size
        self.projection = nn.Conv2d(in_channels, emb_size, 
                                   kernel_size=patch_size, 
                                   stride=patch_size)
        self.cls_token = nn.Parameter(torch.randn(1, 1, emb_size))
        self.positional_embedding = nn.Parameter(
            torch.randn(1, (img_size // patch_size) ** 2 + 1, 
                       emb_size))
    
    def forward(self, x):
        x = self.projection(x)  # [B, D, H/P, W/P]
        x = x.flatten(2).transpose(1, 2)  # [B, N, D]
        batch_size = x.shape[0]
        cls_tokens = self.cls_token.expand(batch_size, -1, -1)
        x = torch.cat((cls_tokens, x), dim=1)  # [B, N+1, D]
        x += self.positional_embedding
        return x
\end{lstlisting}

The patch embedding module performs three operations: (1) convolutional projection of patches into embedding space, (2) addition of a learnable CLS token for classification, and (3) addition of learnable positional embeddings to preserve spatial relationships.

\subsubsection{Multi-Head Self-Attention}

Self-attention allows each patch to attend to all other patches, enabling the model to capture long-range dependencies. The scaled dot-product attention is computed as:

\begin{equation}
\text{Attention}(Q, K, V) = \text{softmax}\left(\frac{QK^T}{\sqrt{d_k}}\right)V
\label{eq:attention}
\end{equation}

where $Q$, $K$, and $V$ are the query, key, and value matrices respectively, and $d_k$ is the dimension of the keys. The scaling factor $\sqrt{d_k}$ prevents the dot products from becoming too large, which would push the softmax into regions with extremely small gradients.

Multi-head attention extends this by computing attention in parallel across $h$ heads, allowing the model to attend to information from different representation subspaces:

\begin{equation}
\text{MultiHead}(Q, K, V) = \text{Concat}(\text{head}_1, \ldots, \text{head}_h)W^O
\label{eq:multihead}
\end{equation}

where $\text{head}_i = \text{Attention}(QW_i^Q, KW_i^K, VW_i^V)$ and $W_i^Q, W_i^K, W_i^V \in \mathbb{R}^{D \times d_k}$, $W^O \in \mathbb{R}^{hd_k \times D}$ are learnable projection matrices.

\begin{lstlisting}[caption={Multi-Head Attention Implementation}]
class MultiHeadAttention(nn.Module):
    def __init__(self, emb_size, num_heads=8, dropout=0.1):
        super().__init__()
        self.num_heads = num_heads
        self.emb_size = emb_size
        self.qkv = nn.Linear(emb_size, emb_size * 3)
        self.projection = nn.Linear(emb_size, emb_size)
    
    def forward(self, x):
        batch_size, tokens, emb_size = x.shape
        qkv = self.qkv(x).reshape(batch_size, tokens, 3, 
                                 self.num_heads, 
                                 emb_size // self.num_heads)
        q, k, v = qkv.permute(2, 0, 3, 1, 4)
        attn = (q @ k.transpose(-2, -1)) * (emb_size ** -0.5)
        attn = attn.softmax(dim=-1)
        x = (attn @ v).transpose(1, 2).reshape(batch_size, 
                                               tokens, emb_size)
        return self.projection(x)
\end{lstlisting}

\subsubsection{Transformer Encoder Layer}

Each encoder layer combines self-attention with a feed-forward network (FFN), using residual connections and layer normalization. The transformer encoder can be expressed as:

\begin{align}
\mathbf{z}'_\ell &= \text{MSA}(\text{LN}(\mathbf{z}_{\ell-1})) + \mathbf{z}_{\ell-1} \label{eq:encoder1} \\
\mathbf{z}_\ell &= \text{MLP}(\text{LN}(\mathbf{z}'_\ell)) + \mathbf{z}'_\ell \label{eq:encoder2}
\end{align}

where $\text{LN}(\cdot)$ denotes Layer Normalization, $\text{MSA}(\cdot)$ is Multi-Head Self-Attention, and $\text{MLP}(\cdot)$ is a two-layer feed-forward network with GELU activation.

\begin{lstlisting}[caption={Transformer Encoder Layer}]
class TransformerEncoderLayer(nn.Module):
    def __init__(self, emb_size, num_heads, 
                 forward_expansion, dropout):
        super().__init__()
        self.layernorm1 = nn.LayerNorm(emb_size)
        self.mha = MultiHeadAttention(emb_size, num_heads)
        self.dropout1 = nn.Dropout(dropout)
        self.layernorm2 = nn.LayerNorm(emb_size)
        self.feed_forward = nn.Sequential(
            nn.Linear(emb_size, forward_expansion * emb_size),
            nn.GELU(),
            nn.Linear(forward_expansion * emb_size, emb_size),
            nn.Dropout(dropout)
        )
    
    def forward(self, x):
        x = x + self.dropout1(self.mha(self.layernorm1(x)))
        x = x + self.feed_forward(self.layernorm2(x))
        return x
\end{lstlisting}

The encoder uses pre-norm architecture (LayerNorm before attention/FFN) with residual connections, enabling stable training of deep networks.

\subsection{Synthesis: How Parts Come Together}

Figure~\ref{fig:architecture} provides a visual overview of how the Vision Transformer components integrate together. The complete Vision Transformer architecture integrates these components as follows:

\begin{lstlisting}[caption={Complete Vision Transformer}]
class VisionTransformer(nn.Module):
    def __init__(self, in_channels=1, patch_size=4, 
                 emb_size=32, img_size=28, num_classes=10, 
                 depth=4, num_heads=4, forward_expansion=2):
        super().__init__()
        self.patch_embedding = PatchEmbedding(...)
        self.transformer_encoders = nn.Sequential(
            *[TransformerEncoderLayer(...) for _ in range(depth)])
        self.layernorm = nn.LayerNorm(emb_size)
        self.mlp_head = nn.Linear(emb_size, num_classes)
    
    def forward(self, x):
        x = self.patch_embedding(x)  # [B, N+1, D]
        x = self.transformer_encoders(x)  # [B, N+1, D]
        x = self.layernorm(x[:, 0])  # Extract CLS token
        return self.mlp_head(x)  # [B, num_classes]
\end{lstlisting}

The forward pass proceeds as: \\(1) Image $\rightarrow$ Patch Embeddings with CLS token (Equation~\ref{eq:patch_embedding}) \\(2) Transformer Encoders process the sequence using self-attention (Equations~\ref{eq:encoder1}--\ref{eq:encoder2}) \\(3) CLS token embedding is extracted \\(4) MLP head produces class predictions. The final classification output is:

\begin{equation}
\mathbf{y} = \text{MLP}(\text{LN}(\mathbf{z}_L^0))
\label{eq:output}
\end{equation}

where $\mathbf{z}_L^0$ is the CLS token representation from the final transformer layer $L$.

\begin{figure}[H]
\centering
\includegraphics[width=0.9\textwidth]{figure1.png}
\caption{Full Vision Transformer pipeline showing patch extraction, embedding, transformer encoder layers, and classification head. The diagram illustrates how input images are divided into patches, embedded with positional information, processed through multiple transformer encoder layers with self-attention, and finally classified using the CLS token.}
\label{fig:architecture}
\end{figure}

\section{Model Training}

\subsection{Training Loop Components}

Our training implementation follows a standard deep learning pipeline with several key components:

\subsubsection{Loss Function}

We use Cross-Entropy Loss, standard for multi-class classification:

\begin{equation}
\mathcal{L}_{\text{CE}} = -\sum_{c=1}^{C} y_c \log(\hat{y}_c)
\label{eq:cross_entropy}
\end{equation}

where $C$ is the number of classes, $y_c$ is the ground truth (one-hot encoded), and $\hat{y}_c$ is the predicted probability for class $c$.

\begin{lstlisting}[caption={Loss Function}]
loss_function = nn.CrossEntropyLoss()
\end{lstlisting}

This loss function computes the negative log-likelihood between predicted class probabilities and true labels, encouraging the model to assign high probability to the correct class.

\subsubsection{Optimizer}

We employ the Adam optimizer with weight decay for regularization:

\begin{lstlisting}[caption={Optimizer Configuration}]
optimizer_adam = optim.Adam(
    vit_model.parameters(), 
    lr=0.0008,
    weight_decay=1e-4,
    betas=(0.9, 0.999)
)
\end{lstlisting}

Adam combines the benefits of AdaGrad (adaptive learning rates per parameter) and RMSProp (moving average of squared gradients), with momentum terms $\beta_1=0.9$ and $\beta_2=0.999$ for first and second moment estimates.

\subsubsection{Learning Rate Scheduling}

We implement a step decay scheduler that reduces the learning rate periodically:

\begin{lstlisting}[caption={Learning Rate Scheduler}]
scheduler = optim.lr_scheduler.StepLR(
    optimizer_adam, 
    step_size=3,
    gamma=0.85
)
\end{lstlisting}

The scheduler multiplies the learning rate by $\gamma=0.85$ every 3 epochs, enabling fine-grained optimization as training progresses. This helps the model converge to better local minima.

\subsubsection{Regularization}

Weight decay ($\lambda = 10^{-4}$) provides L2 regularization, penalizing large parameter values to prevent overfitting. Additionally, we use dropout in the attention and feed-forward layers (set to 0.0 in our miniaturized model for simplicity).

\subsubsection{Training Procedure}

The training loop follows this structure:

\begin{enumerate}
    \item \textbf{Forward Pass:} Compute model predictions for a batch
    \item \textbf{Loss Computation:} Calculate Cross-Entropy loss (Equation~\ref{eq:cross_entropy})
    \item \textbf{Backward Pass:} Compute gradients via backpropagation
    \item \textbf{Optimization Step:} Update parameters using Adam optimizer
    \item \textbf{Metric Tracking:} Update accuracy, F1, precision, recall, AUC metrics
    \item \textbf{Learning Rate Update:} Adjust learning rate via scheduler
\end{enumerate}

During evaluation, we disable gradient computation (\texttt{torch.no\_grad()}) and set the model to evaluation mode (\texttt{model.eval()}), which disables dropout and batch normalization updates.

\subsection{Training Metrics}

We track comprehensive metrics using TorchMetrics:
\begin{itemize}
    \item \textbf{Accuracy:} Overall classification accuracy
    \item \textbf{Top-1 and Top-5 Accuracy:} Standard and top-5 prediction accuracy
    \item \textbf{AUC:} Area Under the ROC Curve for multi-class classification
    \item \textbf{F1 Score:} Harmonic mean of precision and recall
    \item \textbf{Precision and Recall:} Per-class and macro-averaged metrics
\end{itemize}

These metrics provide a comprehensive view of model performance beyond simple accuracy.

\section{Minimal CPU-Ready Working Example (MWE-CPU)}

\subsection{Dataset Selection}

We use a subset of FashionMNIST, a real-world dataset containing $28 \times 28$ grayscale images of 10 clothing categories. The full FashionMNIST dataset contains 60,000 training and 10,000 test samples. For our MWE-CPU, we use:
\begin{itemize}
    \item Training samples: 1,500 (2.5\% of full dataset)
    \item Test samples: 500 (5\% of full dataset)
\end{itemize}

\textbf{Why FashionMNIST for MWE-CPU}\\ \\ \textbf{Dataset Selection Rationale.} We selected FashionMNIST over alternatives 
such as CIFAR-10 for several reasons aligned with our MWE-CPU objectives: 
(1) the grayscale format (single channel) reduces computational overhead and 
simplifies attention pattern visualization, (2) the structured nature of clothing 
items produces clearer patch boundaries compared to natural scenes, (3) lower 
image resolution ($28 \times 28$ vs $32 \times 32$) enables faster CPU training, 
and (4) the well-defined object shapes are well-suited for demonstrating global 
attention mechanisms.


Table~\ref{tab:dataset_comparison} compares FashionMNIST with CIFAR-10 to highlight these advantages.

\begin{table}[H]
\centering
\caption{Dataset Comparison: FashionMNIST vs CIFAR-10}
\label{tab:dataset_comparison}
\begin{tabular}{lcc}
\toprule
\textbf{Feature} & \textbf{FashionMNIST} & \textbf{CIFAR-10} \\
\midrule
Image size & $28 \times 28$ & $32 \times 32$ \\
Channels & 1 (grayscale) & 3 (RGB) \\
Complexity & Lower (structured shapes) & Higher (natural scenes) \\
Noise level & Low & Moderate \\
Interpretability & High (clear boundaries) & Moderate \\
Training speed (CPU) & Fast & Slower \\
\bottomrule
\end{tabular}
\end{table}

This subset size enables fast training on CPU while maintaining meaningful learning signals. FashionMNIST is suitable for ViT as it contains structured visual patterns that benefit from global attention.

\subsection{Model Miniaturization}

To create a CPU-friendly model, we reduce the architecture complexity:

\begin{itemize}
    \item \textbf{Embedding Dimension:} Reduced from 64 to 32
    \item \textbf{Transformer Depth:} Reduced from 7 to 4 layers
    \item \textbf{Attention Heads:} Reduced from 8 to 4 heads
    \item \textbf{Forward Expansion:} Reduced from $4\times$ to $2\times$
\end{itemize}

This results in approximately 36,746 parameters (compared to $\sim$355,000 in a standard configuration), representing an 89\% reduction while maintaining the core transformer architecture.

\subsection{Training Configuration}

Our MWE-CPU training setup is summarized in Table~\ref{tab:config}. These hyperparameters are tuned for stable training on CPU with the miniaturized model.

\begin{table}[H]
\centering
\caption{Model and Training Configuration}
\label{tab:config}
\begin{tabular}{ll}
\toprule
\textbf{Parameter} & \textbf{Value} \\
\midrule
\multicolumn{2}{l}{\textit{Model Architecture}} \\
\midrule
Input channels & 1 (grayscale) \\
Image size & $28 \times 28$ \\
Patch size & $4 \times 4$ \\
Number of patches & 49 ($7 \times 7$) \\
Embedding dimension & 32 \\
Transformer depth (layers) & 4 \\
Number of attention heads & 4 \\
Forward expansion factor & $2\times$ \\
Total parameters & 36,746 \\
\midrule
\multicolumn{2}{l}{\textit{Dataset}} \\
\midrule
Dataset & FashionMNIST (subset) \\
Training samples & 1,500 \\
Test samples & 500 \\
Number of classes & 10 \\
\midrule
\multicolumn{2}{l}{\textit{Training Hyperparameters}} \\
\midrule
Batch size & 64 \\
Epochs & 8 \\
Initial learning rate & 0.0008 \\
Weight decay & $10^{-4}$ \\
Optimizer & Adam ($\beta_1=0.9$, $\beta_2=0.999$) \\
LR scheduler & StepLR (step=3, $\gamma=0.85$) \\
Loss function & Cross-Entropy \\
\midrule
\multicolumn{2}{l}{\textit{Device}} \\
\midrule
Compute device & CPU/MPS/CUDA (auto-detected) \\
\bottomrule
\end{tabular}
\end{table}

\subsection{Training Process and Outcome}

The miniaturized model successfully learns to classify FashionMNIST images. Training completes in minutes on CPU, demonstrating the feasibility of transformer-based vision models on resource-constrained hardware.

\textbf{Training Stability:} Training on grayscale FashionMNIST data demonstrates excellent stability. The monochrome nature reduces input complexity, allowing the model to focus on shape and texture patterns. The learning rate scheduler (StepLR with $\gamma=0.85$ every 3 epochs) ensures smooth convergence without oscillations.

\textbf{Model Size \& Speed Metrics:} Our MWE-CPU implementation achieves:
\begin{itemize}
    \item \textbf{Parameter count:} 36,746 parameters (89\% reduction from standard ViT)
    \item \textbf{Training time per epoch:} $\sim$2--3 minutes on CPU (varies by hardware)
    \item \textbf{Total training time:} $\sim$16--24 minutes for 8 epochs on CPU
    \item \textbf{Memory footprint:} $\sim$150MB during training
\end{itemize}

These metrics validate the MWE-CPU requirement: the model trains efficiently on a standard laptop without GPU acceleration.

Key observations:
\begin{itemize}
    \item The model learns meaningful attention patterns, focusing on discriminative image regions (sleeves, collars, garment shapes)
    \item Despite size reduction, the model captures essential features for classification
    \item Training is stable with appropriate learning rate scheduling
    \item The CLS token effectively aggregates global information from patches
    \item Grayscale input simplifies attention visualization compared to RGB
\end{itemize}

\textbf{Attention Visualization:} Figure~\ref{fig:attention} visualizes the attention patterns learned by the model. The attention visualizations reveal that the model learns to attend to relevant regions of clothing items, such as sleeves, collars, and overall garment shapes. For FashionMNIST, attention patterns differ from CIFAR-10: (1) \textbf{More focused attention} on specific garment parts (e.g., shoe soles, bag handles), (2) \textbf{Clearer boundaries} due to grayscale simplicity, and (3) \textbf{Less spatial diversity} compared to RGB images where attention can vary across color channels. This demonstrates that even a miniaturized ViT can learn interpretable visual representations.

\begin{figure}[H]
    \centering
    \includegraphics[width=0.9\textwidth]{../result/attention_visualization.pdf}
    \caption{Attention visualization showing the model's focus regions for a FashionMNIST sample. The left panel shows the original grayscale image, the middle shows the attention heatmap (warmer colors indicate higher attention), and the right shows the overlay. The model successfully attends to discriminative regions such as garment shapes and key features.}
    \label{fig:attention}
\end{figure}

\textbf{Results Analysis:} The training curves (see Figure~\ref{fig:training_curves}) show consistent learning across epochs, with both training and validation metrics improving over time. The model achieves 59.20\% training accuracy and 61.02\% validation accuracy after 8 epochs, indicating reasonable performance given the small dataset size (1,500 training samples) and miniaturized model architecture. The validation accuracy slightly exceeding training accuracy suggests effective regularization. 

\textbf{Patch Size Impact:} Using $4 \times 4$ patches on $28 \times 28$ images creates 49 patches ($7 \times 7$ grid). This patch size balances between: (1) \textbf{Local detail preservation} (smaller patches capture fine features) and (2) \textbf{Computational efficiency} (fewer patches reduce attention complexity). Larger patches (e.g., $7 \times 7$) would reduce patches to 16 but lose fine-grained details important for distinguishing similar clothing items.

The model achieves strong performance given the small dataset and model size, validating the MWE-CPU approach.

Table~\ref{tab:results} presents the final results after 8 epochs of training, showing all evaluation metrics for both training and validation sets. Figure~\ref{fig:training_curves} visually demonstrates how these metrics change over time throughout the 8 epochs trained.

\begin{figure}[H]
    \centering
    \includegraphics[width=0.9\textwidth]{../result/training_curves.pdf}
    \caption{Training and validation metrics over 8 epochs. The figure shows loss, accuracy, AUC, F1 score, precision, recall, and top-1/top-5 accuracy progression. Both training and validation metrics show consistent improvement, with validation accuracy reaching 61.02\% after 8 epochs.}
    \label{fig:training_curves}
\end{figure}

\begin{table}[H]
\centering
\caption{ViT Results after 8 epochs}
\label{tab:results}
\begin{tabular}{lcc}
\toprule
\textbf{Metric} & \textbf{Train} & \textbf{Validation} \\
\midrule
Loss & 1.1690 & 1.1112 \\
Accuracy & 0.5920 & 0.6102 \\
AUC & 0.9245 & 0.9368 \\
F1 Score & 0.5792 & 0.6138 \\
Precision & 0.5855 & 0.6625 \\
Recall & 0.5920 & 0.6102 \\
Top-1 Accuracy & 0.5920 & 0.6102 \\
Top-5 Accuracy & 0.9757 & 0.9940 \\
\bottomrule
\end{tabular}
\end{table}

\section{Discussion}

\subsection{Limitations}

\textbf{1. Data Efficiency:} Vision Transformers require large amounts of data to perform well. When trained from scratch on small datasets, ViTs typically underperform CNNs due to their lack of inductive biases. Our 61.02\% validation accuracy on 1,500 samples, while demonstrating successful learning, falls short of what CNNs can achieve on similar data scales. This highlights the well-known data hunger of transformer architectures.

\textbf{2. Computational Complexity:} Self-attention has quadratic complexity $O(N^2)$ with respect to sequence length. For our configuration with 49 patches, this is manageable, but scaling to larger images (e.g., $224 \times 224$ with $16 \times 16$ patches yields 196 patches) would significantly increase computation. The attention computation grows from $49^2 = 2,401$ to $196^2 = 38,416$ operations per layer.

\textbf{3. Lack of Inductive Biases:} Unlike CNNs, which have built-in translation equivariance and locality, ViTs must learn these spatial relationships from data. This flexibility comes at the cost of requiring more training examples to achieve comparable performance.

\textbf{4. Fixed Resolution:} Our model uses learnable positional embeddings trained on $28 \times 28$ images. These embeddings do not generalize to different resolutions without interpolation, limiting deployment flexibility.

\textbf{5. Dataset Simplicity:} FashionMNIST's structured, centered objects with uniform backgrounds do not fully challenge the global attention mechanism. More complex datasets with cluttered backgrounds and varying object positions would better demonstrate ViT's strengths.

\subsection{Analysis of Results}

Our model achieved 61.02\% validation accuracy with 59.20\% training accuracy, indicating good generalization without overfitting. The high Top-5 accuracy (99.40\%) suggests the model learns meaningful class relationships even when Top-1 predictions are incorrect. The strong AUC score (0.9368) indicates good discriminative ability across all classes.

The attention visualizations (Figure~\ref{fig:attention}) confirm that the model attends to semantically relevant regions (garment edges, distinctive features) rather than background pixels, validating that the self-attention mechanism successfully identifies discriminative image regions.

\subsection{Future Work}

Based on our findings, we identify several directions for extending this work:

\textbf{1. Baseline Comparison:} A direct comparison with a CNN of similar parameter count (e.g., a small ResNet or custom CNN with $\sim$36K parameters) on the same 1,500-sample subset would quantify the data efficiency gap.

\textbf{2. Data Augmentation:} Implementing augmentations (random crops, rotations, horizontal flips) could improve performance and partially address the data efficiency limitation.

\textbf{3. Knowledge Distillation:} Following DeiT \citep{touvron2021deit}, distilling knowledge from a pre-trained CNN teacher could improve our miniaturized ViT's performance without additional data.

\textbf{4. Self-Supervised Pre-training:} Methods like DINO \citep{caron2021dino} or Masked Autoencoders could enable learning better representations from unlabeled FashionMNIST data before fine-tuning.

\textbf{5. Scaling Study:} Systematically varying model size (embedding dimension, depth, heads) and dataset size would characterize the scaling behavior of miniaturized ViTs.

\section{Conclusion}

This project presented a comprehensive implementation and analysis of the Vision Transformer architecture for image classification. We successfully developed a miniaturized ViT model with 36,746 parameters that trains efficiently on CPU, achieving 61.02\% validation accuracy on a subset of FashionMNIST. Our implementation demonstrates the core principles of ViT: patch-based image representation (Equation~\ref{eq:patch_embedding}), scaled dot-product attention (Equation~\ref{eq:attention}), and transformer encoder layers (Equations~\ref{eq:encoder1}--\ref{eq:encoder2}).

The key contributions of this work include: (1) a fully functional CPU-ready ViT implementation suitable for educational purposes, (2) detailed analysis of how attention mechanisms process grayscale clothing images, and (3) insights into the trade-offs between model size, training efficiency, and classification performance. While ViTs typically require large datasets and computational resources, our MWE-CPU demonstrates that the fundamental architecture can be understood and experimented with on standard hardware.

Future work could explore hybrid CNN-Transformer architectures, self-supervised pre-training on FashionMNIST, and systematic comparison with lightweight CNN baselines on the same data subset.

\bibliographystyle{plainnat}
\bibliography{references}

\end{document}
